\chapter*{Abstract}
\addcontentsline{toc}{chapter}{Abstract}

\selectlanguage{english}

This report presents a method for detecting anomalies in 4G/5G
cellular network traffic, applied to mobile traffic data (Internet,
SMS, calls) from the city of Milan, drawn from the \emph{Telecom
Italia Big Data Challenge} dataset. The proposed approach encodes the
hourly traffic profile of each geographic grid cell as a 24-letter
signature, where each letter represents a traffic level expressed as a
multiple of the cell's hourly median. Each daily signature is then
compared to a reference profile, built by consensus over days of the
same weekday type, using the Damerau-Levenshtein edit distance. The
resulting distances are converted into per-cell z-scores, making it
possible to isolate days with significantly atypical behaviour
($z > 3$). The pipeline is applied independently to the three traffic
channels (Internet, SMS, calls), and detected anomalies are
cross-referenced with a calendar of events (public holidays, AC Milan
and Inter home matches at San Siro). Results show that an anomaly
confirmed simultaneously on at least two channels reaches a 99.2\%
calendar-match rate, compared to 87--98\% for an anomaly detected on a
single channel, confirming the value of a multi-channel approach for
detection reliability. Results can be explored through an interactive
3D dashboard.

\vspace{4mm}
\noindent\textbf{Keywords:} anomaly detection, 4G/5G cellular networks,
mobile traffic, Damerau-Levenshtein distance, z-score, time series,
multi-channel analysis.

\selectlanguage{french}
